%%%%%%%%%%  scan idx 173  |  folio 160  %%%%%%%%%%
\section*{\underline{Exposé II}\\[3pt] \underline{\underline{EXISTENCE DE RÉSOLUTIONS GLOBALES}}}
\begin{center}
\underline{par L. ILLUSIE}
\end{center}

\section*{1. Critères généraux de globalisation}

\textbf{1.0.} Dans ce numéro, on se donne un site $S$, une $S$-catégorie abélienne plate $\mathcal C$ (I 1.1), et une sous-$S$-catégorie additive strictement pleine $\mathcal C_0$ (I 1.1). On suppose $\mathcal C_0$ localement quasi-relevable dans $\mathcal C$ (I 1.2) et localement quasi-stable par noyau d'épimorphisme (I 1.2).

\textbf{Proposition 1.1.} Soit $S\in\Ob S$, et supposons que :
\begin{enumerate}[label=(\roman*)]
\item $\mathcal C_{0S}$ est quasi-relevable dans $\mathcal C_S$ (I 1.2) ($\mathcal C_S$ étant considérée comme catégorie fibrée sur un site ponctuel) ;
\item tout objet de $\mathcal C_S$ qui est de $\mathcal C_0$-type fini (I 1.2) (resp. de $\mathcal C_0$-présentation finie (I 2.8)) est de $\mathcal C_{0S}$-type fini.
\end{enumerate}
Alors :
\begin{enumerate}[label=\alph*)]
\item $\mathcal C_{0S}$ est quasi-stable par noyau d'épimorphisme (I 1.2) ($\mathcal C_{0S}$ étant considérée comme catégorie fibrée sur un site ponctuel).
\item Soient $E\in\Ob \mathrm D^{-}(\mathcal C_S)$ et $n\in\mathbb Z$. Si $E$ est $n$-$\mathcal C_0$-pseudo-cohérent (I 2.3), $E$ est $n$-$\mathcal C_{0S}$-pseudo-cohérent (resp. $(n+1)$-$\mathcal C_{0S}$-pseudo-cohérent). En d'autres termes, on a
\[
\mathrm D^{-}(\mathcal C_S)_{n-\mathcal C_0\text{-coh}}
=
\mathrm D^{-}(\mathcal C_S)_{n-\mathcal C_{0S}\text{-coh}}
\]
(resp.
\[
\mathrm D^{-}(\mathcal C_S)_{n-\mathcal C_0\text{-coh}}
\subset
\mathrm D^{-}(\mathcal C_S)_{(n+1)-\mathcal C_{0S}\text{-coh}}
\subset
\mathrm D^{-}(\mathcal C_S)_{(n+1)-\mathcal C_0\text{-coh}},
\]
inclusions de sous-catégories strictement pleines de $\mathrm D(\mathcal C_S)$, l'inclusion de droite étant d'ailleurs triviale).
\item Pour que $E\in\Ob\mathrm D^{-}(\mathcal C_S)$ soit pseudo-cohérent relativement à $\mathcal C_0$ (I 2.3), il faut et il suffit qu'il le soit relativement à $\mathcal C_{0S}$ ; autrement dit, on a
\[
\mathrm D^{-}(\mathcal C_S)_{\mathcal C_0\text{-coh}}
=
\mathrm D^{-}(\mathcal C_S)_{\mathcal C_{0S}\text{-coh}}.
\]
En outre, le foncteur canonique
\[
\mathrm K^{-}(\mathcal C_{0S})/\mathrm K^{-,\varnothing}(\mathcal C_{0S})
\longrightarrow \mathrm D(\mathcal C_S),
\]
où $\mathrm K^{-,\varnothing}(\mathcal C_{0S})$ désigne la sous-catégorie épaisse de $\mathrm K^{-}(\mathcal C_{0S})$ formée des
com%
